\documentclass{report}
\usepackage{amsfonts}
\newcommand\R{{\mathbb R}}
\newcommand\Q{{\mathbb Q}}
\newcommand\N{{\mathbb N}}
\pagestyle{empty}
\begin{document}
\large
\centerline{\Huge Analisi Matematica II}
\renewcommand{\thefootnote}{ } 
\footnote{\sl Avete 3 ore di tempo. Ogni esercizio
vale sei punti. La sufficienza si ottiene con un punteggio $\geq$
18. Solo le risposte chiaramente giustificate saranno prese in
considerazione.}
\renewcommand{\thefootnote}{\arabic{footnote}} 
\vskip.1cm
\centerline{\Large Prima Prova di Autovalutazione}
\vskip1cm
\begin{enumerate}
\item \label{es1} Si scriva la formula di Taylor per la funzione
\[
f(x)=\sqrt{1-x}.
\]

\item \label{es2} Si scriva la formula di Taylor per la funzione
\[
f(x)=\int_0^1e^{xt^2} dt.
\]
\item Si calcoli il raggio di convergenza per le serie di potenze
trovate negli eserizi \ref{es1} e \ref{es2}.
 
\item Si trovi una successione di numeri razionali la cui somma sia
$\pi$ e la si usi per calcolare $\pi$ con una precisione superiore a
$\frac 1{10}$.
\item Si consideri una funzione $f:[0,1]\to[0,1]$ derivabile in ogni
punto e tale che
\[
\frac 12\leq \frac{f'(x)}{f'(y)}\leq 2 \quad \forall
x,y\in[0,1].
\]
Si dimostri che per ogni coppia di intervalli $I,J\subset [0,1]$
\[
\frac{|I|}{2|J|}\leq\frac{|f(I)|}{|f(J)|}\leq \frac{2|I|}{|J|}
\]
dove, per ogni intervallo $K$, $|K|$ indica la lunghezza di $K$.

\item Si calcoli il limite
\[
\lim_{n\to\infty}\int_{-n}^{n}e^{-x^2}x\,dx
\]
\end{enumerate}
\end{document}
